% Figure: ground-reflection interference. Excess attenuation versus
% frequency for a source and receiver over a hard ground, showing the
% comb of destructive nulls and constructive peaks from the two-path sum.
\begin{figure}[htbp]
\centering
\begin{tikzpicture}
\begin{axis}[
  width=12cm, height=7.5cm,
  xlabel={frequency (Hz)}, ylabel={excess attenuation (dB)},
  xmode=log,
  xmin=50, xmax=2000, ymin=-25, ymax=8,
  axis lines=left,
  domain=50:2000, samples=400,
  legend pos=south west,
  legend cell align=left,
  every axis x label/.style={at={(current axis.right of origin)}, anchor=west},
  every axis y label/.style={at={(current axis.above origin)}, anchor=south},
  grid=major, grid style={enggray!20},
  tick label style={font=\scriptsize},
  label style={font=\small},
  legend style={font=\scriptsize},
]
% Two-path sum over hard ground: direct + reflected with path difference
% delta = 0.34 m, so relative level = 20 log10 | 1 + cos(2 pi f delta / c) |.
% Clamp the argument of the log away from zero so nulls stay finite.
\addplot[engblue, very thick] {20*log10( max(0.056, abs(1 + cos(360*x*0.34/343))) )};
\addlegendentry{hard ground, path difference $0.34$ m}
% the 6 dB constructive-gain reference line
\addplot[enggray, thin, dashed] {6};
\node[enggray, font=\scriptsize, anchor=south east] at (axis cs:1800,6.3)
  {$+6$ dB (in-phase sum)};
% the 0 dB free-field reference
\addplot[enggray, thin] {0};
\end{axis}
\end{tikzpicture}
\caption[Ground-reflection interference and excess attenuation]{Excess
attenuation over free-field level for a source and receiver above a hard
reflecting ground, plotted against frequency for a representative
direct-to-reflected path difference of $0.34$ m. The received field is
the sum of the direct wave and the ground-reflected wave, and their
relative phase turns with frequency, so the level rises toward $+6$ dB
where the two paths arrive in phase and plunges into deep nulls where
they arrive in antiphase. The first null falls where the path difference
is half a wavelength (near $500$ Hz for this geometry), with a comb of
further nulls above it. The pattern is the acoustic ground effect that
shapes every outdoor measurement: moving the microphone height changes
the path difference and slides the comb, which is why standardised
outdoor measurements fix source and receiver heights. Over soft ground
the reflection is weaker and phase-shifted, filling the nulls and adding
a broad low-frequency excess loss instead.}
\label{fig:vol03ch07:ground-interference}
\end{figure}
